% LexGraph: a complete account for a reader who knows nothing about the project.
% Compile with:  tectonic writeup/LexGraph_Explained.tex
\documentclass[12pt,a4paper]{article}

\usepackage[margin=1in]{geometry}
\usepackage{mathptmx}            % Times New Roman text and math
\usepackage[T1]{fontenc}
\usepackage[utf8]{inputenc}
\usepackage{booktabs}
\usepackage{longtable}
\usepackage{array}
\usepackage{amsmath}
\usepackage{amssymb}
\usepackage{enumitem}
\usepackage{xcolor}
\usepackage{tcolorbox}
\usepackage{caption}
\usepackage{titlesec}
\usepackage[hidelinks]{hyperref}

\captionsetup{font=small,labelfont=bf}
\setlist{itemsep=2pt,topsep=3pt,parsep=0pt}
\newcommand{\code}[1]{\texttt{\small #1}}
\newcommand{\pct}[1]{#1\%}
\titleformat{\paragraph}[runin]{\normalfont\bfseries}{}{0pt}{}[.]

\tcbset{
  plainbox/.style={colback=black!3,colframe=black!25,boxrule=0.4pt,
                   left=6pt,right=6pt,top=5pt,bottom=5pt,
                   fonttitle=\bfseries\small,arc=1pt}
}

\title{\textbf{LexGraph}\\[4pt]
\large Predicting Outcomes and Grounding Authorities in\\
Indian Supreme Court Land-Dispute Judgments\\[10pt]
\normalsize A complete account for a reader starting from zero}
\author{CSE 573 --- Semantic Web Mining --- Arizona State University}
\date{1 October 2026}

\begin{document}
\maketitle

\begin{abstract}
\noindent
This document explains, from first principles and assuming no prior knowledge of law or machine
learning, what this project set out to do, what was built, what was measured, and --- in equal
detail --- what turned out not to be true. It is written to be read straight through by someone
who has never seen the repository.

The project began as a system that would read a court judgment, pull structured legal objects out
of it, and reason over them. It became something more useful and less flattering: a carefully
leakage-controlled measurement apparatus which shows that \emph{85\% of the apparent advantage of
the published input construction is disposition language surviving the deletion of the operative
order}, that the court's reasoning itself adds only a small real amount, and that
\emph{converting facts into a symbolic vocabulary destroys most of the predictive signal they
carry}. Both claims are supported by controlled experiments with confidence intervals, and both
contradict what the project originally assumed. Along the way several genuinely positive results
were established --- statute prediction and precedent retrieval both work well above their
baselines, and the system's explanations were shown to be faithful rather than decorative.

Nothing in this document is an estimate or a projection. Every number is read from a results file
produced by a script in this repository, and Section~\ref{sec:status} states for each number
whether it is current or awaiting recomputation.
\end{abstract}

\tableofcontents
\newpage

%==============================================================================
\section{How to read this document}

Section~\ref{sec:problem} explains the problem in ordinary language --- what a court judgment is,
what we are trying to predict, and why anyone would want to. Section~\ref{sec:leakage} explains the
single idea that governs every design decision in the project; if you read only one section, read
that one. Sections~\ref{sec:data} and~\ref{sec:pipeline} describe the data and the machinery.
Section~\ref{sec:results} reports every result. Section~\ref{sec:engineering} is about the
unglamorous work --- the fifteen bugs that would each have produced a confident wrong answer.
Section~\ref{sec:retracted} lists the beliefs the project held and had to give up.
Section~\ref{sec:notestablished} is the honest inventory of what is \emph{not} shown.
Section~\ref{sec:glossary} is a glossary; feel free to jump there at any point.

A reader in a hurry should read Sections~\ref{sec:leakage}, \ref{sec:head1} and \ref{sec:head2},
which contain the whole argument.

%==============================================================================
\section{The problem, in ordinary language}\label{sec:problem}

\subsection{What a judgment is}

When a dispute reaches the Supreme Court of India, the court eventually publishes a document called
a \emph{judgment}. A typical land-dispute judgment in our corpus is around 5,300 words long and has
a predictable shape, though almost never explicit headings:

\begin{enumerate}[label=(\arabic*)]
\item a \textbf{caption}: the case number, the parties' names, the judges, the lawyers;
\item the \textbf{facts}: who owned what, who sold what to whom, when, and what went wrong;
\item the \textbf{procedural history}: which lower courts already heard this and what they decided;
\item the \textbf{arguments}: what each side's counsel submitted;
\item the \textbf{reasoning}: the court's own analysis --- the statutes it reads, the earlier cases
      it follows or distinguishes, and the inferences it draws;
\item the \textbf{order}: the operative disposition. One or two sentences. ``The appeal is
      allowed.'' ``The special leave petition is dismissed with costs.''
\end{enumerate}

Item (6) is the outcome. Items (1)--(4) are, roughly, the inputs that existed before the court
decided. Item (5) is written \emph{afterwards}, by the winner of the internal argument, to justify
a conclusion already reached. That asymmetry is the whole story of this project.

\subsection{The task}

\textbf{Legal judgment prediction} is the task of reading some portion of a case and predicting its
outcome. In our formulation the label is binary:

\begin{itemize}
\item \textbf{WIN} --- the appellant or petitioner succeeded (the appeal was allowed);
\item \textbf{LOSE} --- they did not (dismissed, rejected).
\end{itemize}

Cases that were sent back to a lower court (\code{REMAND}), succeeded in part (\code{PARTIAL}), or
whose disposition our two independent labellers disagreed about (\code{UNRESOLVED}) are held out of
the binary task rather than forced into a class.

\subsection{Why anyone cares}

Two honest motivations and one dishonest one.

\begin{itemize}
\item \textbf{Access to justice.} India's courts carry a backlog measured in tens of millions of
cases. A litigant without money cannot easily learn whether their claim resembles claims that have
succeeded. A system that could say ``here are the rules and the authorities that apply to facts
like yours'' would be useful even if it never predicted an outcome.
\item \textbf{Scientific interest.} If a model can predict outcomes from pre-decision material
alone, that tells us something real about how much of a legal outcome is determined by the facts.
If it cannot, that is equally informative.
\item \textbf{The dishonest one.} ``Our system predicts court decisions with 78\% accuracy'' is a
headline. It is also, as this project demonstrates, substantially an artifact of what the model is
allowed to read.
\end{itemize}

\subsection{What was already published}

The reference point for this project is \textbf{ILDC/CJPE} (Malik et al., \emph{ACL 2021}), which
built a corpus of 34,816 Indian Supreme Court proceedings and reported a best accuracy of
\textbf{78\%}, against \textbf{94\%} for human legal experts.

The ILDC authors are careful and explicit about the obvious trap. They state that the ``end
section(s) directly stating the decision have been deleted from the documents in ILDC since that is
what we aim to predict.'' So ILDC does \emph{not} leak the operative order, and any claim that it
does would be wrong. We checked this and confirmed it.

What ILDC \emph{retains} is item (5) above --- the court's reasoning. The authors say so plainly:
``we consider (along with the facts) the entire case (except the judgment).'' Their best model reads
the \textbf{last 512 tokens} of the document, chosen empirically and justified on the grounds that
``the last parts of case proceedings usually contain the main information about the case and the
rationale behind the judgment.''

That is a legitimate task definition. It is just a different one from the task a litigant faces.
Measuring the difference is this project's central contribution.

\subsection{The two questions}

\begin{tcolorbox}[plainbox,title={The questions this project actually answers}]
\textbf{Q1.} How much of reported legal-judgment-prediction accuracy comes from the court's
reasoning rather than from the facts of the dispute?

\smallskip
\textbf{Q2.} Does converting facts into structured, symbolic legal objects --- the thing an
ontology-driven pipeline is \emph{for} --- help or hurt?
\end{tcolorbox}

The answers, established in Section~\ref{sec:results}, are: \textbf{far less than it appears,
because most of the apparent contribution is not reasoning at all but residual disposition
language}; and \textbf{it hurts, substantially and measurably}.

%==============================================================================
\section{The one idea: leakage}\label{sec:leakage}

\subsection{An analogy}

Suppose you want to predict football results. You build a model, feed it match reports, and report
89\% accuracy. Then someone notices that your match reports include the final score.

Nobody would defend that. But now suppose you carefully delete the score line, and keep everything
else the report says --- including the paragraph that reads ``Liverpool's defence had collapsed by
the hour mark and the result was never in doubt thereafter.'' You have removed the score. You have
not removed the \emph{knowledge} of the score. The report was written by someone who knew how it
ended, and that knowledge is distributed through every sentence.

A court's reasoning section is that paragraph. It is written after the decision, by the judge who
made it, for the purpose of justifying it. The phrase ``we are not persuaded by the appellant's
contention'' is not evidence about the dispute; it is the outcome, restated.

\textbf{Leakage} is the general name for this: information in the model's input that would not be
available at prediction time in the real use case.

\subsection{What must be removed, and why}

The project's masking rule (specified in \code{PROJECT.md} §5.2, implemented in
\code{src/data/mask.py}) keeps only material that existed before the court decided:

\begin{center}
\begin{tabular}{@{}llp{7.6cm}@{}}
\toprule
\textbf{Section} & \textbf{Kept?} & \textbf{Why} \\
\midrule
Caption, parties, judges  & no  & Judge identity is a real predictor but not a legal one; party
                                  names let a model memorise repeat litigants. \\
Facts                     & \textbf{yes} & This is the genuine input. \\
Procedural history        & isolated & Which lower court decided what is \emph{legitimately} known
                                  in advance, but it is also a shortcut, so it is placed in its own
                                  feature group (\code{Q}) and never mixed in silently. \\
Arguments of counsel      & \textbf{yes} & Written by the parties, before the decision. \\
Court's reasoning         & no  & Written after, knowing the answer. \\
Operative order           & no  & This is the label. \\
Headnote (pre-2000)       & no  & An editorial summary added by the law reporter that states both
                                  the outcome and the authorities outright. \\
Page furniture            & no  & ``Indian Kanoon --- \url{http://indiankanoon.org/doc/...}'',
                                  page numbers, running headers. \\
\bottomrule
\end{tabular}
\end{center}

After masking, the average case retains about 2,500 words of the original 5,300.

\subsection{The control that makes everything else interpretable}\label{sec:sensitivity}

Here is the problem with claiming you removed the leakage: how would you know?

The answer is a \textbf{sensitivity control}. We build a deliberately leaky arm --- a model trained
on \emph{only} the text we deleted, the operative order --- and check that it scores near-perfectly.
If it does not, our detector of ``where the order is'' is broken, and the masked arm cannot be
trusted either.

\begin{center}
\begin{tabular}{@{}lrrr@{}}
\toprule
\textbf{Input} & \textbf{Accuracy} & \textbf{Macro-F1} & \textbf{AUROC} \\
\midrule
Majority class (always predict LOSE)        & 0.528 & 0.345 & 0.500 \\
Prior court's disposition only              & 0.562 & 0.562 & 0.585 \\
Masked text (facts + arguments)             & 0.629 & 0.626 & 0.674 \\
Whole unmasked judgment                     & 0.746 & 0.745 & 0.827 \\
\textbf{The removed order only} (control)   & \textbf{0.931} & \textbf{0.931} & \textbf{0.981} \\
\bottomrule
\end{tabular}
\captionof{table}{The leakage probe, \code{forum\_heldout} split, $n_{\text{test}}=902$.
Source: \code{experiments/paper3/leakage\_probe\_forum\_heldout.json}.}
\end{center}

The last row is the point. At AUROC 0.981 the deleted text is almost perfectly predictive, which
confirms two things at once: our order-detector finds the right text, and removing it removes
something that mattered enormously. The gap between 0.827 (whole judgment) and 0.674 (masked) is
the measure of how much the court's own writing was doing.

A second version of the same check uses a large language model instead of a classifier, and agrees:
shown the order alone it reaches 95.7\% accuracy; shown the masked text, 62.0\%
(\code{experiments/paper3/llm\_probe\_forum\_heldout.json}, $n=350$).

\begin{tcolorbox}[plainbox,title={Why this matters more than the headline number}]
Any paper reporting accuracy on this task should report a sensitivity control. Without one, a
reader cannot distinguish ``we removed the leakage'' from ``we believe we removed the leakage''.
This project's masked arm is worth reading precisely because the control next to it is near 1.0.
\end{tcolorbox}

%==============================================================================
\section{The data}\label{sec:data}

\subsection{The corpus}

All 26,688 Indian Supreme Court judgment PDFs from 1950 to 2025 are held locally (6.8\,GB). Of
these, a provided spreadsheet identifies \textbf{6,954} as land or property disputes, and that
subset is the working corpus for everything reported here.

Resolving spreadsheet rows to PDF files sounds trivial and was not. An earlier approach joined on
\emph{normalised titles} and quietly achieved a 0.6\% match rate, because judgment filenames carry
a copy marker after the year (\code{...\_on\_23\_May\_1957\_1}) while titles do not --- so a normaliser
that strips ``trailing digits'' from both removes the \emph{year} from every title, and every pair
then mismatches by exactly that year. The current adapter joins on the exact filename supplied by
the spreadsheet and achieves \textbf{6,954 of 6,954 (100\%)}, with the join rate asserted at
import time (\code{MIN\_PDF\_JOIN\_RATE = 0.99}) so that a regression stops the pipeline instead of
producing a smaller, silently biased dataset.

Text extraction: 6,954 of 6,954 succeeded, with zero failures and zero near-empty results.

\subsection{Where the labels come from}

Nobody hand-labelled 6,954 outcomes. Instead two independent labellers were built and made to
agree:

\begin{enumerate}
\item \textbf{A deterministic rule engine} (\code{src/data/labels.py}) that finds the operative
order by regular expression and classifies it. This is harder than it reads. Indian orders routinely
dispose of several proceedings at once --- ``the appeal is allowed \ldots\ and the writ petition
stands dismissed'' --- and a naive matcher reads the second clause as the outcome. The fix was to
rank proceedings into tiers, so that an \emph{appeal} or \emph{special leave petition} outranks a
consequential \emph{writ petition} or \emph{application}. That change alone reduced
rules-versus-LLM conflicts from 196 to 107.
\item \textbf{A language model} shown the same order window and asked for the same decision, which
must quote the span it relied on.
\end{enumerate}

\textbf{Agreement: 87.0\%} on the 4,317 cases where both decided and the quote verified. Where they
conflict, the case is marked \code{UNRESOLVED} and queued for human review rather than resolved by
coin-flip or by trusting one labeller. There are 562 such cases.

\begin{center}
\begin{tabular}{@{}lrr@{}}
\toprule
\textbf{Label} & \textbf{Cases} & \textbf{Share} \\
\midrule
LOSE        & 2,964 & 42.6\% \\
WIN         & 2,365 & 34.0\% \\
PARTIAL     &   624 &  9.0\% \\
UNRESOLVED  &   562 &  8.1\% \\
UNKNOWN     &   289 &  4.2\% \\
REMAND      &   116 &  1.7\% \\
OTHER       &    34 &  0.5\% \\
\midrule
\textbf{Total} & \textbf{6,954} & \\
\bottomrule
\end{tabular}
\captionof{table}{Merged outcome labels. The binary task uses WIN and LOSE only: 5,329 cases.}
\end{center}

\subsection{How the data is split, and why it is not random}

A random train/test split would be wrong here for two separate reasons.

\paragraph{Reason one: groups} The same dispute often generates several judgments --- an appeal,
a review, a connected matter. If one lands in training and another in test, the model can memorise
the dispute rather than learn anything. So cases are grouped and a group is never split.

\paragraph{Reason two: time} A model that is trained on 2020 cases and tested on 1970 cases is
being allowed to use the future to predict the past. Real deployment never has that.

The project uses two splits, deliberately, because each tests something different:

\begin{itemize}
\item \code{forum\_heldout} --- three forums (courts of origin) are held out entirely. Tests
generalisation to an unseen procedural context. Train 3,350 / dev 455 / test 902.
\item \code{temporal\_2004\_2013} --- train on everything up to 2004, develop on 2004--2013, test
after 2013. Tests generalisation forward in time. Train 3,446 / dev 480 / test 768.
\end{itemize}

Both splits carry a hash of their contents, which every results file records, so a number can
always be traced to the exact split that produced it.

Building the temporal split exposed a conflict between the two requirements above, and a unit test
caught it. Grouping by a case group's \emph{maximum} year put 1970 cases into the ``test after
2013'' bucket; grouping by the \emph{minimum} year put 2025 cases into ``train up to 2004''. Both
are time-travel. There is no assignment that satisfies both constraints for a group that straddles
a boundary, so the 13 straddling groups are dropped and the count is recorded.

%==============================================================================
\section{What was built}\label{sec:pipeline}

The system is a chain of stages, each writing a file the next one reads. Everything is resumable
and everything is deterministic given its inputs.

\subsection{Infrastructure: one door to the model}

All language-model calls go through \code{src/llm/client.py}, which talks to ASU's own hosted model
endpoint (Voyager). Four decisions in that file matter more than they look:

\begin{itemize}
\item \textbf{A content-addressed cache.} Every call is keyed on
\code{sha256(model, messages, parameters)} in SQLite. Re-running an experiment costs nothing --- the
LLM baseline experiment reported later ran at a 100\% cache hit rate over 600 calls. The cache
currently holds 162,230 entries.
\item \textbf{``No answer'' is distinct from ``empty answer''.} A result carries
\code{ok=False} when the call failed, which is \emph{not} the same as a model that replied with
nothing. Conflating the two is how a network outage becomes a scientific finding, which happened
twice on this project before the distinction was enforced (Section~\ref{sec:engineering}).
\item \textbf{A 90-second timeout.} The client library's default is 600 seconds with three
retries, which froze an early batch for over twenty minutes with no output.
\item \textbf{Parallelism.} Extraction runs 96 workers, embeddings 24; throughput went from 10.3 to
38.6 documents per second.
\end{itemize}

\subsection{Segmentation}

Before anything can be extracted, a judgment must be cut into units. The first implementation split
on blank lines and returned \emph{one unit} for 57.7\% of cases, because text extracted from these
PDFs frequently contains no blank lines at all. A third tier --- sliding sentence windows --- was
added. Median units per case went from 1 to 47, and the share of unusable cases from 57.7\% to
2.6\%.

\subsection{Fact extraction, with verbatim grounding}

A model reads each unit and proposes \emph{facts}: atomic factual assertions, each of which must
include a \textbf{verbatim quote} from the source text, plus who asserted it and whether it was
disputed. The quote is then located in the original by
\code{src/grounding.py}; a fact whose quote cannot be found is \textbf{discarded, not corrected}.

This is the project's main defence against fabrication. The current \textbf{grounding rate is
0.92}: 92\% of proposed facts carry a quote that is genuinely present in the judgment.

Getting there required fixing the \emph{verifier}, not the extractor. An early version rejected
42\% of quotes that were in fact correct, for two reasons: PDF text extraction inserts hyphenation
artifacts (\code{plain- tiff's}), and a 20-character minimum rejected short true quotes such as
``Appeal dismissed.'' (15 characters). After repairing hyphenation and lowering the floor, the
verified rate on a hand-checked set went from 58.3\% to 91.7\% --- with 15 of 16 quotes confirmed
present by eye.

\paragraph{Choosing the extraction model} Six models were benchmarked on an identical 20-case set:

\begin{center}
\begin{tabular}{@{}lrrrrr@{}}
\toprule
\textbf{Model} & \textbf{Grounding} & \textbf{Leak} & \textbf{Attribution} & \textbf{Facts/case} & \textbf{s/case} \\
\midrule
\code{llama4-scout-17b}            & 0.92 & 0.03 & 0.27 & 18.4 & \textbf{4.2} \\
\code{llama4-maverick-17b}         & \textbf{0.93} & \textbf{0.01} & \textbf{0.41} & \textbf{23.5} & 28.7 \\
\code{qwen3-30b-a3b-instruct-2507} & 0.92 & 0.03 & 0.40 & 10.5 & 15.7 \\
\code{olmo3-32b-instruct}          & \multicolumn{5}{l}{excluded --- {\raise.17ex\hbox{$\scriptstyle\sim$}}50\,s per call} \\
\code{glm-5-3-flash}               & \multicolumn{5}{l}{excluded --- hangs on long prompts} \\
\code{glm-5-3}                     & \multicolumn{5}{l}{excluded --- too slow at corpus scale} \\
\bottomrule
\end{tabular}
\captionof{table}{Extractor benchmark, 20 cases stratified 1950--2022.
Source: \code{experiments/paper1/extractor\_benchmark.json}.}
\end{center}

\code{maverick} is better on every quality axis and 6.8$\times$ slower. \code{scout} was chosen for
corpus-scale runs; the benchmark is kept so the trade-off is on the record rather than implicit.

\paragraph{Scale} The current fact layer holds \textbf{96,623 facts across 4,702 cases}, median 17
per case.

\subsection{Claims, defences, elements}

A second extraction pass identifies what each party is actually \emph{asking for} --- the claim
(e.g.\ adverse possession, specific performance, an injunction) --- and what is raised against it.
15,377 claim records exist.

The project's ontology (\code{ontology/ontology\_v1.yaml}, converted from a supplied Word document
by \code{src/data/ontology\_convert.py}) defines 15 claims, 66 elements, 6 defence families and 4
remedy tiers. An \emph{element} is a legal requirement a claim must satisfy: adverse possession, for
instance, requires possession that is open, continuous, hostile and for the statutory period.

\subsection{The canonical vocabulary}

Each extracted fact is then \emph{canonicalised}: mapped from free text to a controlled label drawn
from a 345-entry vocabulary generated from the ontology's own named fields. ``The plaintiff had been
cultivating the land since 1962'' becomes something like
\mbox{}\linebreak
\code{possession.continuous\_since\_date}.

This step is the one Section~\ref{sec:head2} shows to be destructive, so it is worth being precise
about what it does and why anyone would want it. The appeal of a controlled vocabulary is that it
makes facts \emph{comparable across cases} --- you can count how often an element is satisfied, mine
patterns, and reason symbolically. The cost is that two facts mapping to the same label become
indistinguishable.

\subsection{Claim families and fact patterns}

With facts as symbols, the specification called for grouping claims into families and then mining
\emph{patterns} --- co-occurring fact sets that predict outcomes within a family. Both were built
(\code{src/cluster/}, \code{src/patterns/}), using HDBSCAN and agglomerative clustering for
families and FP-Growth with closed itemsets for patterns, with chi-square significance testing,
Benjamini--Hochberg correction for multiple comparisons, and bootstrap stability over 20 resamples.

Section~\ref{sec:head2} reports what they found, which is nothing, and why that is itself a result.

\subsection{Grounding facts in law: statutes and precedents}

Two tasks that take facts as input and predict \emph{legal authorities} rather than outcomes.

\paragraph{Statute prediction} Given the facts, predict which statutory provisions apply.
This required building a citation normaliser (\code{src/authorities/normalize.py}) over a registry
of 42 Acts, because real Indian citations arrive in three awkward forms: truncated by line breaks
(the single most frequent ``Act'' in 1,200 sampled judgments is ``Property Act'' with 644 mentions
--- that is the \emph{Transfer of Property Act} with its head cut off); as generic back-references
(``the said Act'', 311+ mentions), resolvable only by looking back to the last Act actually named;
and with implicit Acts (a bare ``Article'' means the Constitution). Coverage over 1,500 judgments:
58,295 citations, 61.4\% resolved to a canonical code.

\paragraph{Precedent retrieval} Given the facts, retrieve the earlier cases the court will cite.
Ten systems were compared, including BM25 (keyword), dense embeddings pooled three different ways,
and reciprocal-rank fusion of both. Every retrieved authority is asserted to pre-date the query
judgment --- checked per result, not merely filtered for.

\subsection{Outcome models and feature groups}

The final predictor consumes the structured pipeline's output as six explicitly separated feature
groups, so that any claim about what drives a prediction can be tested by ablation:

\begin{center}
\begin{tabular}{@{}clr@{}}
\toprule
\textbf{Group} & \textbf{Contents} & \textbf{Features} \\
\midrule
\code{F} & canonical fact atoms            & 1,313 \\
\code{P} & mined fact patterns             & 3 \\
\code{E} & element satisfaction scores     & \textbf{0 (empty --- see B2)} \\
\code{S} & predicted statutes              & 49 \\
\code{R} & retrieved precedents            & 6 \\
\code{C} & case metadata (year, forum, \ldots) & 38 \\
\code{Q} & prior court's disposition       & 6 \\
\bottomrule
\end{tabular}
\captionof{table}{Feature groups, 1,415 features total. \code{Q} is isolated because it is a
procedural shortcut rather than a legal signal.}
\end{center}

\subsection{Traces}

For each prediction the system emits a \emph{trace}: which facts, statutes and precedents it relied
on. Section~\ref{sec:traces} describes the test that establishes these traces are faithful rather
than decorative --- which is not the same as establishing they are correct.

%==============================================================================
\section{Results}\label{sec:results}

\subsection{Headline 1: the published result is mostly leakage}\label{sec:head1}

This is the project's main finding, and it is not the finding the project expected --- an earlier
version of this document claimed the court's \emph{reasoning} was the source of the gap, and
Section~\ref{sec:contamination} reports the measurement that refuted it. Six input constructions were run on \textbf{identical cases,
with an identical split, identical labels and an identical classifier} (TF-IDF features into
logistic regression). The \emph{only} thing that varies between rows is which part of the judgment
the model is allowed to read.

\begin{center}
\begin{tabular}{@{}lrrrr@{}}
\toprule
\textbf{Input construction} & \textbf{Words} & \textbf{Accuracy} & \textbf{Macro-F1} & \textbf{AUROC} \\
\midrule
\code{order\_only} \emph{(sensitivity control)} & 427 & 0.921 & 0.920 & 0.976 \\
\midrule
\code{full} \emph{(whole judgment, ceiling)}    & 5,302 & 0.720 & 0.718 & 0.788 \\
\code{ildc\_style} \emph{(all but the order)}   & 5,051 & 0.680 & 0.678 & 0.734 \\
\code{ildc\_style\_no\_headnote}                & 4,481 & 0.678 & 0.676 & 0.735 \\
\textbf{\code{last\_512\_tokens}} \emph{(ILDC's best input)} & \textbf{511} & \textbf{0.776} & \textbf{0.775} & \textbf{0.845} \\
\midrule
\textbf{\code{masked}} \emph{(facts + arguments only)} & 2,501 & 0.621 & 0.619 & 0.668 \\
\bottomrule
\end{tabular}
\captionof{table}{The reasoning ablation. Same 895 test cases, same 3,324 training cases, same
classifier throughout. Source:
\code{experiments/paper1/reasoning\_ablation\_forum\_heldout.json}.}
\end{center}

Paired bootstrap over the test cases (2,000 resamples), reporting the difference in AUROC with a
95\% confidence interval:

\begin{center}
\begin{tabular}{@{}lrrc@{}}
\toprule
\textbf{Comparison} & \textbf{$\Delta$AUROC} & \textbf{95\% CI} & \textbf{Significant} \\
\midrule
\code{last\_512\_tokens} $-$ \code{masked}        & $+0.1764$ & $[+0.1480, +0.2064]$ & yes \\
\code{ildc\_style\_no\_headnote} $-$ \code{masked} & $+0.0670$ & $[+0.0535, +0.0820]$ & yes \\
\code{ildc\_style} $-$ \code{masked}              & $+0.0660$ & $[+0.0539, +0.0789]$ & yes \\
\code{full} $-$ \code{ildc\_style}                & $+0.0538$ & $[+0.0476, +0.0602]$ & yes \\
\bottomrule
\end{tabular}
\end{center}

\subsubsection*{What this appears to say, and why that reading is wrong}

Read naively, the table says the court's reasoning is worth a great deal. Take \code{order\_only} as
the ceiling on what is knowable here (AUROC 0.976, so $0.476$ of signal is available above chance);
then the facts and arguments appear to carry $(0.668-0.5)/0.476 = 35\%$ of it, the whole
order-removed document $49\%$, and the last 512 tokens $73\%$.

That reading has one weak point, and it turns out to be fatal to it. \textbf{The final 512 tokens sit
exactly where the operative order used to be.} If the boundary between ``reasoning'' and ``order'' is
even slightly imperfect, that arm --- and no other --- is contaminated. The next section measures
whether it is.

\subsubsection*{Three objections, and what answers each}

This is the kind of result that invites immediate objections. Each was tested; the second one won.

\paragraph{``Your masked arm is worse simply because it has less text''} No --- the evidence runs
hard the other way. The \emph{best} non-leaky arm, \code{last\_512\_tokens}, reads \textbf{511
words}; the masked arm reads \textbf{2,501}. The winning construction has roughly one-fifth the text
and beats the masked arm by 0.176 AUROC. Volume is not the mechanism; \emph{which} text is.

\paragraph{``The last 512 tokens must still contain the order''} \textbf{This objection is correct},
and Section~\ref{sec:contamination} is the measurement that establishes it. An earlier version of
this document argued the objection was unlikely, on two indirect grounds: that the order detector is
independently validated (the \code{order\_only} arm reaches 0.976, so it finds the order rather than
missing it), and that \code{ildc\_style} beats \code{masked} without any positional advantage. Both
of those facts are true and neither settles the question, because the order detector can mark the
boundary correctly and still leave the sentence \emph{before} it in place. That is exactly what
happens.

\paragraph{``This is just your corpus''} Correct, and this is the single most important
qualification on the result. Everything above is measured on \emph{our} 6,954 land-dispute cases
with \emph{our} masking, as a comparable construction of ILDC's task --- not on ILDC itself. The
replication on ILDC's own released text and labels is written
(\code{src/eval/ildc\_ablation.py}, four arms including a positional control) and blocked only on
dataset licence acceptance, which is the user's to give and which this project deliberately did not
route around.

\subsubsection*{One number worth keeping}

The \code{last\_512\_tokens} arm reaches \textbf{accuracy 0.776}; ILDC's published best is
\textbf{78\%}. A plain TF-IDF model, on a different and much smaller corpus, lands within half a
percentage point of a published transformer result once it reads the same window the published model
reads. That is worth keeping --- but after the next section it is no longer a statement about
reasoning. It is a statement about how much can be recovered from a window that both setups leave
contaminated.

\subsection{Where the advantage actually comes from}\label{sec:contamination}

The claim to be tested is that \code{last\_512\_tokens} wins because the final window retains
disposition language, not because the court's reasoning is informative. Three measurements settle
it, on 842 test cases (60 of the 902 were dropped for being shorter than 2,048 words in the
order-removed text, which the setback arms require). \textbf{Every arm here is exactly 512 words by
construction}, so volume cannot explain any difference between them.

\subsubsection*{1. The setback curve}

Take four non-overlapping 512-word windows, stepping backwards from the end of the order-removed
document. The two hypotheses predict different shapes, and both predictions were written into the
module before it was run:

\begin{itemize}
\item \textbf{Residual leakage} --- contamination lives next to the boundary and cannot reach back
512 words, so the curve should \emph{fall off a cliff} after the first window.
\item \textbf{Reasoning} --- a court's analysis occupies the last several thousand words, so signal
should \emph{decay gradually} with distance, with every window still above \code{masked}.
\end{itemize}

\begin{center}
\begin{tabular}{@{}llrrrc@{}}
\toprule
\textbf{Arm} & \textbf{Window} & \textbf{AUROC} & \textbf{$\Delta$ vs \code{masked}} & \textbf{95\% CI} & \textbf{Sig.} \\
\midrule
\code{masked} & facts + arguments (2,616 w) & 0.656 & --- & --- & --- \\
\midrule
\code{w0} & $[-512{:}]$      & \textbf{0.838} & $+0.1814$ & $[+0.1501, +0.2134]$ & \textbf{yes} \\
\code{w1} & $[-1024{:}-512]$ & 0.669 & $+0.0128$ & $[-0.0158, +0.0410]$ & \textbf{no} \\
\code{w2} & $[-1536{:}-1024]$ & 0.647 & $-0.0094$ & $[-0.0358, +0.0154]$ & no \\
\code{w3} & $[-2048{:}-1536]$ & 0.638 & $-0.0183$ & $[-0.0457, +0.0070]$ & no \\
\bottomrule
\end{tabular}
\captionof{table}{The setback curve. Source:
\code{experiments/paper1/leakage\_setback\_forum\_heldout.json}.}
\end{center}

This is unambiguously the leakage shape: a cliff, not a decay. \textbf{512 words of the court's
reasoning, taken from anywhere other than the very end, carry no more outcome signal than the facts
do.}

\subsubsection*{2. A cue audit, against the right base rate}

Why that shape? Run the disposition patterns from \code{src/data/labels.py} --- the \emph{same}
regular expressions that produced this project's outcome labels --- over each window and count the
cases containing an explicit disposition statement.

A raw count would be uninterpretable. Masked text is full of sentences like ``the High Court
dismissed the suit'', which recite what a \emph{lower} court did and leak nothing about this court's
decision; the patterns match those too. So \code{masked} supplies the base rate, and only an
\emph{excess} over it is evidence of contamination.

\begin{center}
\begin{tabular}{@{}lrrl@{}}
\toprule
\textbf{Arm} & \textbf{Cases with a cue} & \textbf{Rate} & \textbf{Excess over base} \\
\midrule
\code{masked} & 227 / 842 & 0.270 & \emph{--- base rate} \\
\code{w0} & 353 / 842 & \textbf{0.419} & $\mathbf{+0.150}$ \\
\code{w1} & 75 / 842 & 0.089 & $-0.180$ \\
\code{w2} & 70 / 842 & 0.083 & $-0.186$ \\
\code{w3} & 62 / 842 & 0.074 & $-0.196$ \\
\bottomrule
\end{tabular}
\end{center}

Only the final window is \emph{enriched} in disposition language. Every setback window sits well
\emph{below} the base rate --- cleaner than \code{masked} --- and scores the same as it.

\paragraph{The mechanism, visible in the text} Indian judgments announce the conclusion one sentence
before the operative order:

\begin{quote}
``For the foregoing reasons we find no merit in the appeal. \emph{The appeal is dismissed with
costs.}''
\end{quote}

The order detector marks the boundary at the second sentence, which is \textbf{correct} --- that is
the operative order. The first sentence stays, and it states the outcome just as plainly.

\subsubsection*{3. How deep the contamination reaches}

The setback curve answers the question at 512-word granularity. Sweeping the same 512-word window
backwards in fine steps locates the boundary precisely --- and AUROC and cue rate fall together,
which is what a single shared cause looks like:

\begin{center}
\begin{tabular}{@{}rrrrc@{}}
\toprule
\textbf{Offset} & \textbf{AUROC} & \textbf{Cue rate} & \textbf{$\Delta$ vs \code{masked}} & \textbf{Sig.} \\
\midrule
0 words   & 0.838 & 0.419 & $+0.1814$ & yes \\
32 words  & 0.780 & 0.296 & $+0.1240$ & yes \\
64 words  & 0.719 & 0.208 & $+0.0625$ & yes \\
\textbf{128 words} & 0.684 & 0.131 & $+0.0278$ & \textbf{no} \\
256 words & 0.676 & 0.099 & $+0.0200$ & no \\
512 words & 0.669 & 0.089 & $+0.0128$ & no \\
\bottomrule
\end{tabular}
\captionof{table}{Contamination depth. The advantage is gone once the window is set back
\textbf{128 words}.}
\end{center}

\textbf{Deleting the operative order leaves roughly 100 words of outcome-telegraphing text behind,
and that residue accounts for essentially all of the last-512 arm's advantage.}

Two subsidiary checks point the same way and bound how much a cue-detector alone can do. Deleting
every cue-bearing \emph{sentence} from \code{w0} drops it from 0.838 to 0.775; evaluating the same
model only on the 489 cases where no cue was detected gives 0.742, against \code{masked}'s 0.646 on
those same cases. Both remain above \code{masked}, so sentence-level scrubbing does \emph{not} fully
clean the window --- there is outcome-telegraphing language the patterns miss (``we find no merit in
the submission'', ``the conclusion is inescapable''). This is exactly why the setback curve is the
primary evidence: it never relies on detecting a cue at all.

\subsubsection*{So does the reasoning add anything? Yes, but little}

One loose end remains. The whole-document arm \code{ildc\_style} \emph{contains} the contaminated
tail, so its advantage might be the same residue diluted across 5,000 words. Cutting the tail off
answers it:

\begin{center}
\begin{tabular}{@{}lrrc@{}}
\toprule
\textbf{Comparison} & \textbf{$\Delta$AUROC} & \textbf{95\% CI} & \textbf{Sig.} \\
\midrule
\code{w0} (last 512) $-$ \code{masked} & $+0.1814$ & $[+0.1501, +0.2134]$ & yes \\
\code{ildc\_style} $-$ \code{masked} & $+0.0588$ & $[+0.0458, +0.0722]$ & yes \\
\code{ildc\_style} $-$ \code{ildc\_minus128} \emph{(the tail alone)} & $+0.0322$ & $[+0.0273, +0.0376]$ & yes \\
\midrule
\textbf{\code{ildc\_minus128}} $-$ \code{masked} \emph{(reasoning, cleaned)} & $\mathbf{+0.0266}$ & $\mathbf{[+0.0143, +0.0391]}$ & \textbf{yes} \\
\code{ildc\_minus512} $-$ \code{masked} & $+0.0215$ & $[+0.0096, +0.0334]$ & yes \\
\bottomrule
\end{tabular}
\end{center}

With the final 128 words removed, the court's \emph{entire} reasoning still beats facts-and-arguments
by $+0.027$ AUROC, and the interval excludes zero. So the reasoning does carry real signal --- it is
simply far smaller than it appears. Of \code{ildc\_style}'s $+0.059$ advantage, more than half
($+0.032$) is the contaminated tail.

\begin{tcolorbox}[plainbox,title={The corrected headline}]
The input construction used by the published state of the art appears to beat a facts-only input by
$\mathbf{+0.181}$ AUROC. \textbf{About 85\% of that is an artifact of imperfect order removal.}
Deleting the operative order leaves roughly 100 words of outcome-telegraphing text in place; a
512-word window set back past that residue scores no better than the facts at all
($+0.013$, CI $[-0.016, +0.041]$).

\smallskip
With the contaminated tail removed, the court's reasoning adds $\mathbf{+0.027}$ AUROC
($\mathbf{[+0.014, +0.039]}$) --- real, statistically significant, and \textbf{one-seventh} of what
it appears to add.
\end{tcolorbox}

\paragraph{Why this is a better result than the one it replaces} The earlier claim --- ``the court's
reasoning carries more signal than the facts'' --- was a claim about legal text. This one is a claim
about \emph{benchmark construction}, it is more precisely quantified, and it is actionable: a
benchmark built this way needs roughly 128 words of margin past the detected order boundary, not
zero. It also explains a puzzle the earlier framing left open, namely why a 511-word window would
beat a 2,501-word one so decisively when it supposedly contained only a subset of the same kind of
reasoning.

\paragraph{One caveat on scope} These arms run on 842 cases rather than the 895 of
Section~\ref{sec:head1}, because the setback windows need documents of at least 2,048 words. The two
tables are therefore not directly comparable row by row --- \code{masked} is 0.656 here and 0.668
there --- but every comparison \emph{within} each table is paired on identical cases, which is what
the confidence intervals are computed over.

\subsection{Headline 2: turning facts into symbols destroys the signal}\label{sec:head2}

The second question was whether the structured, ontology-driven representation helps. To answer it
cleanly, the project built a \textbf{representation ladder}: the same classifier, the same split, the
same labels and the same 881 test cases at every rung, with \emph{only the representation} changing.

\begin{center}
\begin{tabular}{@{}llrrr@{}}
\toprule
\textbf{Rung} & \textbf{Representation} & \textbf{Features} & \textbf{Macro-F1} & \textbf{AUROC} \\
\midrule
0 & majority class                 & ---    & 0.346 & --- \\
1 & masked judgment text           & 50,000 & 0.608 & \textbf{0.657} \\
2 & extracted fact text            & 41,142 & 0.586 & 0.614 \\
3 & canonical fact atoms           & 1,313  & 0.505 & \textbf{0.517} \\
\bottomrule
\end{tabular}
\captionof{table}{The representation ladder, \code{forum\_heldout}, $n_{\text{test}}=881$.
Source: \code{experiments/paper1/representation\_ladder\_forum\_heldout.json}.}
\end{center}

Read down the AUROC column. Extraction costs \textbf{0.043}. Canonicalisation --- converting those
extracted facts into the controlled vocabulary --- costs a further \textbf{0.097}, landing at 0.517,
which is essentially chance. The pipeline's symbolic output carries almost no information about the
outcome, while the raw text it was derived from carries a measurable amount.

\subsubsection*{Five independent lines of evidence for the same conclusion}

A single ladder could be a bug. Five separate tests were run and all agree.

\begin{enumerate}
\item \textbf{The ladder itself}: $0.657 \rightarrow 0.614 \rightarrow 0.517$.

\item \textbf{Is it the \emph{ontology's} vocabulary that is bad?} No. An A/B test compared the
authored ontology vocabulary against a vocabulary \emph{induced} from the corpus at equal
granularity. Raw text beats both decisively, and the two vocabularies are statistically
indistinguishable from each other:
\begin{center}
\begin{tabular}{@{}lrrr@{}}
\toprule
\textbf{Comparison} & \textbf{$\Delta$AUROC} & \textbf{95\% CI} & \textbf{$P$} \\
\midrule
text $-$ induced   & $+0.1138$ & $[+0.0732, +0.1545]$ & 0.000 \\
text $-$ ontology  & $+0.0887$ & $[+0.0430, +0.1352]$ & 0.000 \\
induced $-$ ontology & $-0.0251$ & $[-0.0711, +0.0200]$ & \textbf{0.846} \\
\bottomrule
\end{tabular}
\end{center}
So the loss is not a defect of \emph{this} vocabulary. It is a property of discretisation.

\item \textbf{Would more labels help?} No. A sweep from 345 to 1,000 vocabulary entries does not
recover the gap, and a frequency-threshold sweep sits at chance at every threshold.

\item \textbf{Does the loss happen where the vocabulary fails, or where it succeeds?} This is the
sharpest test (\code{src/eval/vocab\_diagnosis.py}). Facts were divided into those the vocabulary
\emph{can} name and those it cannot, and the text-versus-atom comparison was run within each group:
\begin{center}
\begin{tabular}{@{}lrrr@{}}
\toprule
\textbf{Group} & \textbf{Text AUROC} & \textbf{Atom AUROC} & \textbf{Cost of naming} \\
\midrule
A. facts the vocabulary \emph{can} name    & 0.544 & 0.457 & \textbf{0.087} \\
B. facts it \emph{cannot} name             & 0.526 & 0.524 & 0.002 \\
\bottomrule
\end{tabular}
\end{center}
The cost is concentrated \emph{precisely where the vocabulary works}. That rules out the comfortable
explanation (``coverage is just incomplete; fix the vocabulary'') and supports the uncomfortable one:
successfully assigning a fact to a legal category is what discards the information. Note the small
sample ($n_{\text{test}}=143$ and $127$) --- group A's interval spans zero at $P=0.083$, so this
line is suggestive rather than conclusive on its own.

\item \textbf{A pre-registered prediction.} Before running it, the project predicted that an LLM
given canonical facts (\code{LLM-CF}) would do \emph{worse} than the same LLM given raw text
(\code{LLM-0}). It did: macro-F1 0.469 versus 0.491.
\end{enumerate}

\subsubsection*{The downstream consequence: §8 finds nothing, correctly}

If canonical atoms carry no signal, then pattern mining over canonical atoms must find nothing. It
does, and the machinery was built well enough to say so rather than to manufacture a result:

\begin{center}
\begin{tabular}{@{}lr@{}}
\toprule
Closed itemsets mined                                   & 498 \\
Stable across $\geq$80\% of 20 bootstrap resamples      & 311 \\
Patterns statistically tested                            & 351 \\
Pass uncorrected $p<0.05$                                & 10 \\
\emph{Expected to pass by chance alone}                  & \emph{17.6} \\
\textbf{Survive Benjamini--Hochberg correction}          & \textbf{0} \\
\bottomrule
\end{tabular}
\end{center}

Ten patterns passed an uncorrected test --- fewer than the 17.6 chance alone predicts. Without the
multiple-comparisons correction, those ten would have shipped as a ``pattern library''. This is the
clearest example in the project of a control earning its place.

\paragraph{The claim taxonomy also fails to validate} Clustering 7,441 claims produced six clusters
with a silhouette score of \textbf{0.0245} and 71.5\% of points classed as outliers by HDBSCAN. A
silhouette near zero means the clusters are not separated. When the clusters were merged into
families with a minimum size, one family (\code{PrivateTitlePossession}) absorbed 2,365 cases and
nearly every other label. The taxonomy is not supported by the data, and is reported as such rather
than presented as a finding.

\subsection{Where structured facts \emph{do} pay off}\label{sec:positives}

The two results above are negative. They are not the whole story: on the tasks that predict
\emph{legal authorities} rather than outcomes, facts work well.

\subsubsection*{Statute prediction}

Given only the facts, predict the statutory provisions that apply. 103 labels,
$n_{\text{test}}=796$.

The first version of this experiment produced micro-F1 0.447 and looked like a strong result. It was
an artifact, and the thing that caught it was a \textbf{copy control}: a system that does nothing but
copy provisions already mentioned verbatim in the input text scores \textbf{0.713} --- far higher.
The ``prediction'' was largely string matching.

The honest version removes from the target set every provision the input text mentions, leaving only
provisions the model must genuinely infer. 46 labels, $n_{\text{test}}=702$:

\begin{center}
\begin{tabular}{@{}lrrr@{}}
\toprule
\textbf{System} & \textbf{Micro-F1} & \textbf{P@1} & \textbf{R@10} \\
\midrule
Global top-5 prior                     & 0.069 & 0.134 & 0.418 \\
Claim-family prior \emph{(strongest baseline)} & 0.093 & 0.173 & 0.503 \\
Copy from text \emph{(control, now at floor)} & 0.000 & 0.081 & 0.385 \\
\midrule
\textbf{Masked text}                   & \textbf{0.173} & \textbf{0.236} & \textbf{0.660} \\
Extracted fact text                    & 0.159 & 0.206 & 0.581 \\
Canonical atoms                        & 0.109 & 0.161 & 0.464 \\
\bottomrule
\end{tabular}
\captionof{table}{Statute prediction on \emph{novel} labels only --- provisions not mentioned in the
input. Source: \code{experiments/paper2/statute\_prediction\_forum\_heldout\_labels\_novel.json}.}
\end{center}

\textbf{Micro-F1 0.173 against the strongest baseline's 0.093 --- 1.9$\times$}, with the copy control
driven to zero by construction. This is a real positive result, and the same ordering appears again:
raw text $>$ extracted facts $>$ canonical atoms.

\subsubsection*{Precedent retrieval}

Given the facts, retrieve the earlier cases the court will actually cite. 150 queries. An important
honesty constraint applies: only \textbf{1,035 of 2,113} cited authorities (49.0\%) are themselves in
this corpus, because Supreme Court judgments cite High Courts, the Privy Council and English
decisions. Recall is reported against the \emph{reachable} set with the raw figure stated beside it.

\begin{center}
\begin{tabular}{@{}lrrrr@{}}
\toprule
\textbf{System} & \textbf{R@10} & \textbf{R@50} & \textbf{MRR} & \textbf{nDCG@10} \\
\midrule
Mention baseline                       & 0.022 & 0.112 & 0.037 & 0.014 \\
Popularity baseline \emph{(control)}   & 0.032 & 0.103 & 0.072 & 0.032 \\
Canonical atom overlap                 & 0.021 & 0.058 & 0.046 & 0.021 \\
Dense, mean-pooled                     & 0.053 & 0.130 & 0.114 & 0.054 \\
Dense, max-pooled                      & 0.039 & 0.104 & 0.097 & 0.044 \\
BM25 (keyword)                         & 0.207 & 0.350 & \textbf{0.431} & 0.225 \\
Dense over masked text                 & 0.229 & 0.390 & 0.380 & 0.221 \\
\textbf{RRF(BM25 + dense masked text)} & \textbf{0.238} & \textbf{0.453} & 0.428 & \textbf{0.243} \\
\bottomrule
\end{tabular}
\captionof{table}{Precedent retrieval, 150 queries, recall against 1,035 reachable authorities.
Source: \code{experiments/paper2/precedent\_retrieval\_forum\_heldout.json}.}
\end{center}

The fused system reaches \textbf{R@10 0.238 and MRR 0.431 against a popularity control of 0.032 and
0.072} --- roughly \textbf{7.4$\times$} the control on R@10. Because the control is popularity (always
return the most-cited cases), beating it by that margin establishes the system is responding to the
facts rather than exploiting citation skew.

Two sub-results worth noting. Canonical atom overlap (0.021) is \emph{below} the popularity control
--- the symbolic representation fails on this task too. And max-pooling the dense representation made
it worse (0.039 vs 0.053), which was a prediction this project made and had to retract: the deficit
is in the facts, not in the pooling.

\subsection{Outcome prediction hits a ceiling, and the ceiling looks real}\label{sec:ceiling}

Twenty-four outcome models were run on 881 test cases across 1,415 features. Selected rows:

\begin{center}
\begin{tabular}{@{}lrrrr@{}}
\toprule
\textbf{Model} & \textbf{Accuracy} & \textbf{Macro-F1} & \textbf{95\% CI} & \textbf{ECE} \\
\midrule
Majority (global)              & 0.530 & 0.346 & $[0.332, 0.360]$ & --- \\
\textbf{Majority (per decade)} & 0.571 & \textbf{0.557} & $[0.527, 0.589]$ & --- \\
\midrule
\code{lr F} (facts only)       & 0.501 & 0.499 & $[0.464, 0.530]$ & 0.270 \\
\code{lr F+P+E+S+R+C}          & 0.513 & 0.513 & $[0.479, 0.546]$ & 0.285 \\
\code{lr S only} (statutes)    & 0.599 & 0.599 & $[0.567, 0.629]$ & \textbf{0.041} \\
\code{lr Q only} (prior court) & 0.525 & 0.509 & $[0.474, 0.540]$ & 0.034 \\
\midrule
\code{gbm C only} (metadata)   & 0.583 & 0.582 & $[0.550, 0.613]$ & 0.091 \\
\textbf{\code{gbm F+P+E+S+R+C}} & \textbf{0.598} & \textbf{0.597} & $[0.566, 0.630]$ & 0.107 \\
\bottomrule
\end{tabular}
\captionof{table}{Outcome models, $n_{\text{test}}=881$. ECE is expected calibration error --- how
far the model's stated confidence is from its observed accuracy; lower is better.
Source: \code{experiments/paper3/outcome\_forum\_heldout.json}.}
\end{center}

Three uncomfortable observations:

\begin{enumerate}
\item \textbf{Every confidence interval spans about 0.11 and they all overlap.} The best structured
model (0.597) and a per-decade majority baseline (0.557) are not distinguishable at this sample size.
Reporting the former as a result without the latter beside it would be misleading.
\item \textbf{Adding facts to metadata does not help.} \code{gbm C only} --- year and forum, no facts
at all --- reaches 0.582. Facts alone reach 0.517. McNemar's exact test confirms the
metadata-and-authority groups genuinely beat the fact group ($P=0.0008$ for logistic regression,
$P=0.004$ for gradient boosting), which is a real difference in the wrong direction for the
project's original thesis.
\item \textbf{Calibration and accuracy come apart.} \code{lr S only} has the best calibration in the
table (ECE 0.041) at competitive accuracy, while the best-accuracy model is poorly calibrated
(0.107). For a system meant to tell a litigant how confident it is, the calibrated model may be the
more useful one.
\end{enumerate}

\subsubsection*{Large language models do no better}

Five LLM configurations on the same 300-case subset, same labels:

\begin{center}
\begin{tabular}{@{}lrrrr@{}}
\toprule
\textbf{System} & \textbf{Accuracy} & \textbf{Macro-F1} & \textbf{AUROC} & \textbf{ECE} \\
\midrule
\code{LLM-0} (zero-shot)            & 0.577 & 0.491 & 0.579 & 0.265 \\
\code{LLM-FS} (few-shot, $k=4$)     & 0.543 & 0.522 & 0.522 & 0.117 \\
\code{LLM-CoT} (chain of thought)   & 0.583 & 0.582 & 0.610 & 0.207 \\
\code{LLM-CF} (canonical facts)     & 0.560 & 0.469 & 0.545 & 0.062 \\
\code{LLM-CF+A} (canonical + authorities) & 0.560 & 0.503 & 0.535 & 0.046 \\
\midrule
Structured \code{gbm} (same cases)  & 0.603 & 0.599 & 0.628 & 0.109 \\
TF-IDF on masked text               & 0.597 & 0.593 & \textbf{0.645} & 0.045 \\
Majority (per decade)               & 0.587 & 0.571 & --- & --- \\
\bottomrule
\end{tabular}
\captionof{table}{Head-to-head, $n=300$, \code{llama4-scout-17b}.
Sources: \code{experiments/paper3/llm\_baselines\_forum\_heldout.json},
\code{head2head\_forum\_heldout.json}.}
\end{center}

Chain-of-thought's apparent advantage over zero-shot is not significant (McNemar $P=0.926$). Nothing
in the table clears a per-decade majority baseline by a margin the sample size can resolve. The
ceiling is not an artifact of choosing a weak model family.

\subsubsection*{Error analysis suggests the ceiling is intrinsic}

354 errors on 881 cases were categorised, with a \textbf{base-rate control} that turns out to be
decisive. For each diagnostic signal, compare the error rate among cases where the signal fires to
the overall error rate of 0.402. The ratio is the \emph{lift}; a lift of 1.0 means the signal tells
you nothing.

\begin{center}
\begin{tabular}{@{}lrrr@{}}
\toprule
\textbf{Signal} & \textbf{Errors} & \textbf{$P$(error $\mid$ flag)} & \textbf{Lift} \\
\midrule
\code{family\_unassigned} & 39  & 0.448 & 1.12 \\
\code{statute\_mismatch}  & 175 & 0.398 & 0.99 \\
\code{discretionary}      & 67  & 0.372 & 0.93 \\
\code{extraction\_thin}   & 25  & 0.352 & 0.88 \\
\code{atoms\_only\_generic} & 1 & 0.250 & 0.62 \\
\bottomrule
\end{tabular}
\captionof{table}{Error-signal diagnosticity against a base error rate of 0.402.
Source: \code{experiments/paper3/error\_analysis\_forum\_heldout.json}.}
\end{center}

\textbf{Every lift is approximately 1.0.} Not one of the pipeline's own quality signals predicts
where it will fail. Had these been reported as raw counts --- ``175 errors involve a statute
mismatch'' --- they would have read as a diagnosis and an obvious work plan. Against the base rate
they are noise.

The second test is stronger. If the model's errors were caused by \emph{pipeline} defects, they
should differ from the errors of a completely different system. On the 300 shared cases the
structured model made 119 errors and \code{LLM-CoT} made 125; under independence one expects
$119 \times 125 / 300 = 49.6$ cases wrong for both. The observed number is \textbf{46}. The two
systems' errors are statistically independent --- they are not failing on the same cases, and
therefore not failing for a shared, fixable reason. The remaining difficulty looks intrinsic to
predicting outcomes from pre-decision facts in this corpus.

\subsection{The explanations are faithful}\label{sec:traces}

A system that outputs a prediction plus a list of ``the facts I relied on'' invites a simple
question: did it rely on them, or is the list decoration?

The test is a \textbf{deletion test}. Delete the facts the trace cites, re-run the model, and measure
how much the predicted probability moves. Compare against deleting the same number of random facts,
and against deleting the facts the trace ranked \emph{least} important.

One design decision matters here. The deletion is applied to the \textbf{facts}, and features are
recomputed from scratch --- not to the feature vector directly. Deleting features would be circular:
of course removing a feature the model weighted heavily changes the output. Deleting the underlying
facts tests whether the trace points at the right \emph{evidence}.

\begin{center}
\begin{tabular}{@{}lrrc@{}}
\toprule
\textbf{Facts deleted} & \textbf{Mean $\Delta$probability} & \textbf{vs cited} & \textbf{95\% CI} \\
\midrule
\textbf{Cited by the trace} & \textbf{0.0650} & --- & --- \\
Random                      & 0.0209 & $+0.0441$ & $[+0.0245, +0.0659]$ \\
Inverse-ranked              & 0.0180 & $+0.0470$ & $[+0.0291, +0.0666]$ \\
\bottomrule
\end{tabular}
\captionof{table}{Trace faithfulness, 45 cases. Cited facts move the prediction 3.1$\times$ as much
as random facts, winning on 77.8\% of cases.
Source: \code{experiments/paper3/trace\_faithfulness\_forum\_heldout\_gbm.json}.}
\end{center}

\textbf{Verdict: faithful.} The facts the trace names are the facts the model used.

\begin{tcolorbox}[plainbox,title={What faithfulness is not}]
Faithful does not mean \emph{correct}. This test establishes that the explanation describes the
model's actual computation. It says nothing about whether that computation is legally sound. Given
that the same model predicts outcomes near chance (Section~\ref{sec:ceiling}), a faithful
explanation of a near-chance prediction is an honest account of a weak inference --- which is the
right thing to produce, but should not be oversold.
\end{tcolorbox}

%==============================================================================
\section{What it took: the bugs that mattered}\label{sec:engineering}

A reader evaluating this project should know that most of the effort went not into building the
pipeline but into establishing that its numbers meant what they appeared to mean. Fifteen defects
were found and fixed; each would have produced a confident, publishable, wrong answer.

{\small
\begin{longtable}{@{}p{4.6cm}p{5.2cm}p{4.4cm}@{}}
\toprule
\textbf{Defect} & \textbf{What it would have caused} & \textbf{How it was caught} \\
\midrule
\endfirsthead
\toprule
\textbf{Defect} & \textbf{What it would have caused} & \textbf{How it was caught} \\
\midrule
\endhead
Segmentation returned one unit for 57.7\% of cases (no blank lines in PDF text)
 & Extraction silently operating on 2.6\% of its intended input
 & Median-units sanity check \\
\addlinespace
Quote verifier rejected 42\% of \emph{correct} quotes (PDF hyphenation; 20-char floor)
 & Grounding rate understated; good extractions discarded
 & Hand audit of 16 quotes \\
\addlinespace
Consequential dismissals read as the disposition
 & 196 cases labelled with the wrong outcome
 & Auditing rules-vs-LLM conflicts \\
\addlinespace
Group-whole and time-respect split constraints in direct conflict
 & Training on the future to predict the past
 & \textbf{A unit test} (\code{tests/test\_stage1.py}) \\
\addlinespace
Fact attribution read only \code{asserted\_by}, 84\% \code{court\_narrative}
 & Attribution coverage reported as 17\% instead of 96\%
 & Inspecting the field distribution \\
\addlinespace
Canonicaliser's self-reported ``NEW'' rate of 0.7\%
 & A vocabulary declared complete while 16.7\% of facts were unnameable
 & Checking the answers instead of the flag \\
\addlinespace
Batching the canonicaliser with a shared label menu
 & Label fidelity 34\% $\rightarrow$ 22.7\%, uninvestigated
 & Per-batch-size sweep (4/8/16) \\
\addlinespace
\code{canonicalize\_one} silently dead after a prompt-ID change
 & The \emph{fallback} path returning \code{None} for everything
 & Reading the validation path \\
\addlinespace
Driver ran two pipeline stages concurrently; only the last exit code checked
 & A failed stage advancing the chain as if it had succeeded
 & Noticing a 26-minute orphaned process \\
\addlinespace
macOS idle sleep mistaken for a hang (stalls of 78, 194, 232 min)
 & Three runs killed and restarted for no reason
 & \code{pmset -g log} \\
\addlinespace
Burden annotation: 92\% one class, one value never used
 & Legal metadata that was a constant, read downstream as fact
 & A degeneracy gate that \emph{refuses} to save \\
\addlinespace
Order-window detector swallowed whole short documents (59 cases)
 & A third of all unusable cases, cause unknown
 & Unusable-case audit \\
\addlinespace
Caption detector ate argument text (\code{For the Appellant}; page footers)
 & \textbf{545 of 5,255 units} --- {\raise.17ex\hbox{$\scriptstyle\sim$}}1.2M words corpus-wide --- deleted as ``caption''
 & Reading what was being dropped \\
\addlinespace
Loop-variable shadowing in the metrics dump
 & Metrics overwritten by a loop counter
 & Implausible output value \\
\addlinespace
Per-item LLM loop with no retries (twice)
 & A network outage reported as ``the corpus does not state these elements''
 & \code{ok=False} $\neq$ empty answer \\
\bottomrule
\end{longtable}
}

\subsection{Four transferable lessons}

\begin{enumerate}
\item \textbf{A control is worth more than a method.} The copy control killed a statute result
(0.447 $\rightarrow$ honest 0.173); the popularity control validated a precedent result; the
base-rate control dissolved an entire error taxonomy; Benjamini--Hochberg prevented a pattern library
from shipping. In every case the control changed the conclusion.
\item \textbf{Distinguish ``no answer'' from ``the answer is nothing''.} This single confusion
produced two separate false findings on this project. A loop that makes one call per item needs
retries and must count failures separately from empty results.
\item \textbf{Assert your joins.} The PDF join and the fact-extraction QA gate both assert their
success rate and abort below it. A pipeline that silently processes 12\% of its input is worse than
one that crashes.
\item \textbf{Elapsed time is not progress, and 0\% CPU is not a hang.} Three runs were killed
before anyone checked the power log.
\end{enumerate}

%==============================================================================
\section{Claims this project held and gave up}\label{sec:retracted}

Research projects usually report what survived. This one kept a record of what did not, because the
retractions are more informative than the confirmations.

\begin{center}
\small
\begin{tabular}{@{}p{5.4cm}p{8.6cm}@{}}
\toprule
\textbf{Belief} & \textbf{What the measurement said} \\
\midrule
``The last-512 window is clean, so its $+0.176$ advantage shows the court's \emph{reasoning} carries
the signal.''
 & \textbf{Wrong, and this document said so before it was tested.} The window is cue-enriched
   ($+0.150$ over the base rate) and the advantage vanishes once it is set back 128 words. Roughly
   85\% of it is residual disposition language; the reasoning's real contribution is $+0.027$.
   See Section~\ref{sec:contamination}. \\
\addlinespace
``Extraction preserves the signal in the text.''
 & At $n=881$ it loses 0.043 AUROC. Canonicalisation loses a further 0.097. \\
\addlinespace
``The ontology's vocabulary is the problem; an induced one will do better.''
 & Indistinguishable: induced $-$ ontology $= -0.025$, $P=0.846$. The loss is discretisation
   itself. \\
\addlinespace
``The procedural shortcut (prior court's disposition) dominates everything.''
 & Weak. \code{Q only} reaches AUROC 0.498 at $n=881$ --- chance. \\
\addlinespace
``Max-pooling will fix the dense retrieval arm.''
 & Made it worse: R@10 0.039 vs 0.053 mean-pooled. The deficit is in the facts. \\
\addlinespace
``ILDC leaks the operative order.''
 & It does not. The authors explicitly delete it, and we verified the claim. What ILDC retains is
   the \emph{reasoning}, which is a different and defensible choice. \\
\addlinespace
``Recovering the 1.2M words of wrongly-deleted argument text will improve fact attribution.''
 & Flat: 0.20 $\rightarrow$ 0.19. Absolute attributed facts rose {\raise.17ex\hbox{$\scriptstyle\sim$}}15\% (15,899 $\rightarrow$ 18,358) but the
   \emph{rate} did not move. The cause of the 20\% attribution ceiling is still unknown. \\
\addlinespace
``Swapping in a better grounding model will fix the missing \code{Fact} and \code{Issue} concepts.''
 & Tested across five models. Rules learned under one model fire on nothing under another, and the
   recovered rules were single-example artifacts that generalised to zero real sentences. \\
\addlinespace
``A bigger model will produce usable element-burden metadata.''
 & Two models failed the same degeneracy gate; the larger, more instruction-disciplined one failed
   too. \\
\bottomrule
\end{tabular}
\end{center}

The last two illustrate a pattern worth naming: on this project, \emph{model capability was almost
never the binding constraint}. Instruction discipline, evaluation design and data coverage were.

%==============================================================================
\section{What is \emph{not} established}\label{sec:notestablished}

This section exists so that no reader has to discover these limits for themselves.

\subsection{Open blockers}

\paragraph{B2 --- element burden metadata does not exist} The ontology requires every element to
record who bears the burden of proof, to what standard, and when it shifts. It supplies this for
\textbf{0 of 66} elements. Two models were asked and both produced degenerate output (92\% and 87\%
assigned to a single class; one standard never used at all, including where it unambiguously applies).
A gate refuses to save such annotations, so feature group \code{E} is \textbf{empty} and every result
labelled \code{F+P+E+...} contains no \code{E}. The fix is a few hours of expert hand annotation,
not a larger model.

\paragraph{B3 --- the vocabulary cannot be frozen} Measured over 79,496 canonical facts,
\textbf{20.3\% are out of vocabulary} against a 5\% threshold for freezing. The seed vocabulary of
345 subcategories covers roughly three quarters of what real judgments state. Note that
Section~\ref{sec:head2} shows improving this coverage would not rescue the predictive result --- the
cost of naming is concentrated where the vocabulary already succeeds.

\paragraph{B4 --- there is no gold annotation set} The specification asks for at least 150 cases
annotated by a legal reviewer with inter-annotator agreement measured on a 30-case overlap. The
\textbf{150 cases are prepared} (\code{Data/gold/gold\_tasks.jsonl}) and an annotation tool is built
(\code{src/annotate/app.py}), but no reviewer has annotated them. Consequently \textbf{extraction
precision and recall against gold cannot be computed}. Four proxies are in use, and it is worth being
exact about what each does and does not establish:

\begin{center}
\begin{tabular}{@{}lp{9.2cm}@{}}
\toprule
\textbf{Proxy} & \textbf{What it establishes} \\
\midrule
Quote grounding rate (0.92)   & Bounds \emph{fabrication}. Says nothing about recall. \\
Legal-statement filter        & Bounds one \emph{precision} failure mode. \\
Rules-vs-LLM agreement (87.0\%) & Estimates \emph{outcome-label} noise, not fact quality. \\
The representation ladder     & Measures \emph{downstream utility} --- what we actually care about ---
                                but cannot attribute a loss to any specific extraction error. \\
\bottomrule
\end{tabular}
\end{center}

This is the single biggest gap for a reviewer, and it is blocked on human time, not compute.

\subsection{Other open items}

\begin{itemize}
\item \textbf{The ILDC replication is not done.} Everything in Section~\ref{sec:head1} is measured on
our corpus as a comparable construction of ILDC's task. The module that runs the same ablation on
ILDC's own released text and labels is written and tested; the dataset is gated behind a licence whose
acceptance belongs to the user, and the module refuses to route around the gate. Until it runs, the
headline is a strong internally-controlled result, not a replication.
\item \textbf{What the cue scrub cannot separate.} The residual-leakage audit is done
(Section~\ref{sec:contamination}) and found contamination reaching about 128 words. But
\code{w0\_scrubbed} and \code{w0\_cue\_free} both remain above \code{masked}, so some
outcome-telegraphing language in the final window escapes the patterns. Whether that remainder is
undetected leakage or genuine summative reasoning is \textbf{not} resolved --- the setback curve
sidesteps the question rather than answering it.
\item \textbf{Human verification of outcome labels.} The specification's check of 200 labels is
outstanding; the 562 \code{UNRESOLVED} conflicts are where it should start.
\item \textbf{The attribution ceiling.} 19--20\% of facts are attributed to a party. The hypothesis
that missing argument text caused this was tested and refuted. The real cause is unknown.
\item \textbf{Every outcome-model confidence interval spans about 0.11.} At $n=881$ the structured
model and a per-decade majority baseline cannot be separated. Larger $n$ or a different task is
required to say more.
\item \textbf{§8 as specified is mis-specified for this corpus.} Pattern mining over canonical atoms
cannot work when canonical atoms carry no signal. This is reported as a finding rather than quietly
dropped, but it does mean a specified deliverable is vacuous rather than met.
\end{itemize}

%==============================================================================
\section{Status of every number in this document}\label{sec:status}

A fact-layer defect was fixed on 30 September 2026: the caption detector had been deleting argument
text (row 13 of the table in Section~\ref{sec:engineering}). Fixing it changed the fact layer
substantially --- 79,496 facts $\rightarrow$ \textbf{96,623}; 4,581 cases $\rightarrow$ 4,702; median
13 per case $\rightarrow$ 17. Everything downstream of facts is therefore being recomputed, and this
section says plainly which numbers are current.

\begin{center}
\begin{tabular}{@{}llp{6.4cm}@{}}
\toprule
\textbf{Result} & \textbf{Status} & \textbf{Why} \\
\midrule
Leakage probe (§\ref{sec:sensitivity})      & \textbf{current} & TF-IDF over text; independent of the fact layer \\
Reasoning ablation (§\ref{sec:head1})       & \textbf{current} & Re-run after the fix \\
Setback / contamination (§\ref{sec:contamination}) & \textbf{current} & New, 1 Oct 2026; TF-IDF over text \\
Corpus, labels, splits (§\ref{sec:data})    & \textbf{current} & Upstream of facts \\
Extractor benchmark (§\ref{sec:pipeline})   & \textbf{current} & Fixed 20-case set \\
Fact counts, grounding (§\ref{sec:pipeline}) & \textbf{current} & Post-fix: 96,623 facts, 0 dropped, QA 4,026/4,026 spans exact \\
\midrule
Representation ladder rung 1                & \textbf{current} & Text-based rung \\
Representation ladder rungs 2--3            & \emph{recomputing} & Derived from facts \\
Vocabulary A/B and diagnosis                & \emph{recomputing} & Derived from facts \\
Statute and precedent results               & \emph{recomputing} & Fact-derived arms; text arms unaffected \\
Outcome models, LLM baselines               & \emph{recomputing} & Feature groups \code{F},\code{P} change \\
Patterns, claim families                    & \emph{recomputing} & Derived from facts \\
Trace faithfulness, error analysis          & \emph{recomputing} & Derived from the outcome model \\
\bottomrule
\end{tabular}
\end{center}

\paragraph{What the re-run cannot change} The two headline findings do not depend on the fact layer.
Section~\ref{sec:head1} is TF-IDF over text throughout. Section~\ref{sec:head2}'s \emph{direction} is
structural --- a 1,313-feature discretisation of a 50,000-feature text signal cannot overtake it by
gaining 21\% more facts --- though the exact magnitudes will move. The honest expectation is that the
ladder gaps narrow slightly and the conclusion stands; if it does not, that will be reported.

At the time of writing, canonicalisation of the new fact layer is 77\% complete
(74,485 of 96,623).

%==============================================================================
\section{Where the project stands}\label{sec:standing}

\subsection{What exists and works}

\begin{itemize}
\item A corpus of 6,954 land-dispute judgments, 100\% resolved to source PDFs, fully text-extracted.
\item Dual-labeller outcome labels at 87.0\% agreement, with conflicts preserved rather than resolved.
\item A masking implementation with a sensitivity control at AUROC 0.981 --- the apparatus that makes
every other number interpretable.
\item Two principled splits, hashed, with group and time constraints enforced and tested.
\item A fact extraction layer of 96,623 verbatim-grounded facts at a 0.92 grounding rate.
\item A statute normaliser over 42 Acts handling truncation, back-references and implicit Acts.
\item Working statute prediction (1.9$\times$ baseline) and precedent retrieval (7.4$\times$ control).
\item A faithfulness-tested trace mechanism.
\item A reproducible driver (\code{run\_downstream.sh}) with per-stage abort-on-failure and a QA gate.
\end{itemize}

\subsection{What the paper would say}

The defensible contribution is not a system that predicts outcomes well --- it does not. It is a
measurement:

\begin{tcolorbox}[plainbox,title={The contribution}]
\textbf{1.} On Indian Supreme Court land disputes, the input construction used by the published
state of the art beats a facts-only input by $+0.181$ AUROC --- and \textbf{about 85\% of that gap
is an artifact of imperfect order removal}, not the court's reasoning. Deleting the operative order
leaves roughly \textbf{128 words} of outcome-telegraphing text behind; a 512-word window set back
past it scores no better than the facts ($+0.013$, CI $[-0.016, +0.041]$). With the tail removed,
the reasoning's real contribution is $+0.027$ AUROC, CI $[+0.014, +0.039]$. The practical
consequence for anyone building such a benchmark is a concrete margin, not a caution.

\smallskip
\textbf{2.} Projecting facts into a symbolic legal vocabulary \textbf{costs 0.097 AUROC}, which is
most of the signal that survives extraction. This is a property of discretisation, not of the
particular vocabulary: an induced vocabulary of equal granularity is statistically indistinguishable
($P=0.846$), more labels do not help, and the cost is concentrated precisely where the vocabulary
succeeds in naming a fact.

\smallskip
\textbf{3.} Facts \emph{are} predictive of \textbf{legal authorities} --- statutes at 1.9$\times$ the
strongest baseline, precedents at 7.4$\times$ a popularity control --- which is where an
ontology-driven pipeline should be pointed instead.
\end{tcolorbox}

Result 3 is what rescues the project from being purely negative, and it is also the professor's
original framing: facts in, rules and authorities out.

\subsection{What comes next, in priority order}

\begin{enumerate}
\item \textbf{Replicate on ILDC itself.} The single step that would move this from a
strong internally-controlled result to a replication. Written and waiting on a licence.
\item \textbf{Get the 150 gold cases annotated.} Removes the largest reviewer objection: there is
currently no precision or recall figure for extraction, only proxies.
\item \textbf{Refresh all fact-derived results} once the driver chain completes, and update the
reports.
\item \textbf{Hand-annotate the 66 element burdens} (B2), which unblocks feature group \code{E} and
the element-level reasoning the ontology was built for.
\end{enumerate}

%==============================================================================
\section{Glossary}\label{sec:glossary}

\begin{description}[leftmargin=0pt,style=nextline,itemsep=3pt]
\item[AUROC] Area under the receiver-operating-characteristic curve. The probability that the model
scores a randomly chosen WIN case above a randomly chosen LOSE case. 0.5 is chance; 1.0 is perfect.
Preferred here over accuracy because it does not depend on where the decision threshold is set.
\item[Macro-F1] The F1 score computed per class and then averaged, so that doing well on the
majority class alone cannot carry the number. A model that always predicts LOSE gets 53\% accuracy
but only 0.346 macro-F1.
\item[ECE (expected calibration error)] How far a model's stated confidence sits from its observed
accuracy. A model that says ``70\% confident'' and is right 70\% of the time has ECE near 0.
\item[Paired bootstrap] To compare two systems, resample the test cases with replacement 2,000
times, recompute both scores on each resample, and look at the distribution of the \emph{difference}.
If the 95\% interval excludes zero, the difference is unlikely to be sampling noise. ``Paired'' means
both systems are scored on the same resample, which removes case difficulty as a confound.
\item[McNemar's test] An exact test for whether two classifiers differ, which looks only at the cases
where they disagree. Appropriate here because the two systems are evaluated on the same cases.
\item[Benjamini--Hochberg correction] When 351 hypotheses are tested at $p<0.05$, roughly 17 will pass
by chance. This procedure adjusts the threshold to control the proportion of false discoveries among
those reported. Applying it took this project's pattern library from 10 patterns to zero.
\item[Lift (base-rate control)] The error rate among flagged cases divided by the overall error rate.
A lift of 1.0 means the flag carries no information, however large its raw count.
\item[TF-IDF] Term frequency $\times$ inverse document frequency: a representation of a document as
weighted word counts, where words common across the whole corpus are down-weighted. Simple, fast,
and in this project repeatedly competitive with much larger models.
\item[BM25] A keyword-ranking function used by search engines. On this project's precedent-retrieval
task it beat dense embeddings on long queries.
\item[Dense retrieval] Representing documents as vectors from a neural embedding model and retrieving
by vector similarity --- as opposed to keyword matching.
\item[RRF (reciprocal rank fusion)] A way of combining two ranked lists by summing the reciprocals of
each item's rank in each list. Used here to fuse BM25 with dense retrieval, which beat both.
\item[Closed itemset] In pattern mining, a set of co-occurring items that has no superset occurring
equally often. Mining only closed itemsets avoids reporting thousands of redundant sub-patterns.
\item[Leakage] Information in the model's input that would not be available when the prediction is
actually needed. The organising concern of this entire project; see Section~\ref{sec:leakage}.
\item[Sensitivity control] A deliberately leaky arm, included to prove the leakage detector works. If
a model trained on the text you \emph{removed} does not score highly, you did not remove what you
thought you did.
\item[Copy control] A baseline that does nothing but copy answers already present verbatim in the
input. If it beats your system, your system is doing string matching. This one invalidated a statute
result.
\item[Popularity control] A baseline that always returns the most frequently cited items. Beating it
shows a retrieval system is responding to the query rather than exploiting skew in the data.
\item[Grounding] Requiring a model to quote its source verbatim, then verifying the quote appears in
the original. An unverifiable quote means the fact is discarded.
\item[Trace] The list of facts, statutes and precedents a prediction was based on.
\item[Faithfulness] Whether a trace describes the model's actual computation. Distinct from, and much
weaker than, correctness.
\item[IRAC] Issue, Rule, Application, Conclusion --- the standard structure of a legal analysis, and
the output format the system's reports follow.
\item[Adverse possession] A legal doctrine by which someone who openly occupies land that is not
theirs, continuously and against the owner's interest, for a statutory period, may acquire title.
Prominent in this corpus and used throughout as an example.
\item[SLP (special leave petition)] The discretionary route by which most cases reach the Indian
Supreme Court. Its disposition is the outcome label for a large share of this corpus.
\end{description}

%==============================================================================
\appendix
\section{How to reproduce everything}

\begin{enumerate}
\item \textbf{Environment.} \code{.venv/bin/pip install -r requirements.txt}. Create \code{.env} at
the project root with \code{VOYAGER\_API\_KEY=<key>}, obtained from
\code{voyager.rc.asu.edu}~$\rightarrow$~LLM Access~$\rightarrow$~Create Key. No VPN is required; the
endpoint resolves to public addresses.
\item \textbf{Run the chain.} \code{./run\_downstream.sh}. It re-executes itself under
\code{caffeinate -dimsu} so macOS idle sleep cannot stall it, aborts on any stage's non-zero exit,
and runs a QA gate after extraction. Every stage is resumable.
\item \textbf{Headline experiments individually:}
\begin{quote}\footnotesize\ttfamily\raggedright
python -m src.eval.leakage\_probe --split forum\_heldout\\
python -m src.eval.reasoning\_ablation --split forum\_heldout\\
python -m src.eval.leakage\_setback --split forum\_heldout\\
python -m src.eval.representation\_ladder --split forum\_heldout\\
python -m src.eval.vocab\_ab --split forum\_heldout\\
python -m src.authorities.statute\_predict --split forum\_heldout\\
\hspace*{2em}--rebuild-targets\\
python -m src.authorities.precedent\_retrieve --split forum\_heldout\\
python -m src.predict.models --split forum\_heldout\\
python -m src.trace.evaluate --split forum\_heldout --model gbm
\end{quote}
\item \textbf{Every number in this document} is read from a file under \code{experiments/}. The
figure generator (\code{src/report/figures.py}) reads the same files and \emph{raises} rather than
substituting a default if one is missing, so a stale figure cannot be produced silently.
\end{enumerate}

\section{Where things live}

\begin{center}
\small
\begin{tabular}{@{}lp{9.4cm}@{}}
\toprule
\textbf{Path} & \textbf{Contents} \\
\midrule
\code{PROJECT.md}       & The specification. Section numbers (§5.2, §6.2, §10.1) refer to it. \\
\code{WORKLOG.md}       & Chronological record of decisions, dead ends and bugs with reasoning. \\
\code{reports/}         & Stage acceptance against the specification: \code{M0}--\code{M5},
                          \code{BLOCKERS.md}, \code{ENV\_NOTES.md}. \\
\code{src/data/}        & Adapter, labels, masking, segmentation, splits, ontology conversion. \\
\code{src/extract/}     & Fact and claim extraction, canonicalisation, vocabulary. \\
\code{src/eval/}        & The experiments: leakage probe, reasoning ablation, representation
                          ladder, vocabulary A/B and diagnosis, error analysis, ILDC replication. \\
\code{src/authorities/} & Statute normalisation and prediction, precedent retrieval. \\
\code{src/predict/}     & Feature construction, outcome models, LLM baselines. \\
\code{src/trace/}       & Trace construction and the faithfulness deletion test. \\
\code{src/llm/client.py} & The single door to the model endpoint: cache, timeout, retries. \\
\code{src/grounding.py} & Verbatim quote location with hyphenation repair. \\
\code{experiments/}     & Every result file. Nothing in this document is sourced elsewhere. \\
\code{Data/interim/}    & Derived artifacts. Large ones are gitignored and rebuilt deterministically. \\
\code{Data/gold/}       & The 150 prepared gold cases and the 66-element burden sheet (B2, B4). \\
\code{run\_downstream.sh} & The driver. \\
\bottomrule
\end{tabular}
\end{center}

\vfill
\begin{center}
\small\emph{Every figure in this document was read from a results file in \code{experiments/}.
Where a number is awaiting recomputation, Section~\ref{sec:status} says so.}
\end{center}

\end{document}
